% !TeX root = ../PIFB2.tex
\section{Discussion: What the Variational Principle Contributes}\label{sec:discussion}

The principal gain is a consistent account of how local descriptions are compared, revised and retained. A normalized joint yields an evidence bound. Selective channels make the comparison physically interpretable as partial access. A joint statistical bundle avoids treating beliefs and models as unrelated ontological substances. Conditional expectation and prediction generate cross-scale inference once a retention and dynamics are supplied.

The principle does not eliminate the need for a model. Changing a prior, a measurement interface, a noise law or a mobility changes the question being answered. Some of those ingredients can themselves be inferred or learned under a larger model. Enlarging the model shifts the boundary of what is explained; it does not make the boundary disappear. The productive task is to derive more structure while tracking that shift explicitly.

\subsection{Persistence, geometry and mechanics}
Precision has several distinguishable roles. It can enter a restoring Hessian, the Fisher metric or an observational likelihood. First-order natural-gradient dynamics turn the metric into dissipative resistance. A separately postulated kinetic metric defines a different dynamics. The exact memory construction of Appendix~\ref{app:examples} demonstrates persistence without conflating these roles.

Likewise, Fisher geometry and Lorentz geometry answer different questions. The former measures statistical sensitivity; the latter arises conditionally from a determinant cone and its transformations. No contradiction is resolved by pretending the positive Fisher form becomes indefinite. The open physical task is to derive the structure that relates these two constructions to records of motion and propagation.

\subsection{Participatory interpretation}
The mathematical framework is compatible with a participatory reading because records, selective communication and retained collective descriptions affect subsequent predictions. Its stronger claim would be that such dynamics generate the stable, recursively organized objects through which a physical world is described. The four-cycle selection theorem addresses one part of that claim in one family. Exact parent laws and memory identify the next requirements.

Interpretation does not follow uniquely from the equations. A realist may read them as models of situated access to an underlying world; a structural interpretation may emphasize invariant relations among records. The framework's empirical content must come from specified predictions, not from selecting a philosophical vocabulary.

